% ============================================================================
% Fichier  : partie2/P02_C07_decomposer_espace_audit.tex
% Titre    : Décomposer l'Espace d'Audit : Gabarit Canonique et Fiche d'Assertion
% Auteur   : MINKA MI NGUIDJOÏ Thierry Emmanuel
% Version  : 1.0 -- octobre 2026
% ============================================================================
% Sources : article A0 « Vers l'auditologie » (HAL hal-05701007),
% proposition de nécessité d'une décomposition hiérarchique et thématique T2
% (gabarit canonique à huit niveaux, conditions de validité) ; article A4
% « Hierarchical Decomposition in Auditology » (en préparation, non cité en
% bibliographie) : contraintes d'exhaustivité, de non-redondance et de borne
% cognitive, conditions de richesse, de cohérence sémantique et
% d'hétérogénéité bornée, adaptations. Règle d'atomicité et attributs de la
% fiche d'assertion : cours Audit_SI de l'auteur, mis en cohérence le
% 30/10/2026 (décision de l'auteur : option (a)).
% ============================================================================

\chapter[Décomposer l'Espace d'Audit]{Décomposer l'Espace d'Audit : Gabarit Canonique et Fiche d'Assertion}
\label{chap:decomposition}
\chaptermark{Décomposer l'espace d'audit}

\epigraph{%
    « Ce qu'on ne peut pas tenir en tête,\\
    on ne peut pas l'auditer~: on le découpe. »%
}{Principe de la méthode}

\niveauChapitre{Intermédiaire}{%
    Maîtrise des Chapitres~\ref{chap:controle-interne} (contrôles)
    et~\ref{chap:planifier} (plan d'audit)%
}

\begin{boxobjectifs}
\begin{itemize}
    \item Justifier la décomposition de l'espace d'audit par les limites de
          la mémoire de travail.
    \item Décomposer un domaine selon le gabarit canonique à huit niveaux~:
          domaine, sous-domaine, processus, activité, contrôle, test, preuve,
          trace.
    \item Rédiger la fiche d'assertion d'un contrôle~: objectif, critère,
          question, risque, conséquence.
    \item Appliquer la règle d'atomicité aux objectifs, aux critères et aux
          questions.
    \item Vérifier qu'une décomposition est exhaustive, non redondante et
          respecte la borne cognitive.
    \item Reconnaître les domaines où le gabarit doit être adapté.
\end{itemize}
\end{boxobjectifs}

\bigskip

% ============================================================================
\section*{Question directrice : pourquoi la grille de 45 questions a-t-elle échoué~?}
% ============================================================================

Le prologue a donné deux raisons à l'échec de la première mission de
LiZbiZ~: une grille binaire qui interdisait de dire « je ne sais pas », et
une grille reprise d'une autre entreprise, sans critères propres. Il y en a
une troisième. Les neuf questions du domaine~D mélangeaient des objets de
niveaux différents~: une question portait sur la politique d'accès, une
autre sur un journal technique, une troisième combinait deux conditions.
Personne ne pouvait dire si le domaine avait été couvert en entier.

La seconde mission a commencé par \emph{décomposer} le domaine~D en
$5 \times 5$ éléments~\cite{minka2026auditologie}. Ce chapitre enseigne
cette opération. Il en donne la raison (section~\ref{sec:c07-pourquoi}), le
gabarit (section~\ref{sec:c07-gabarit}), la fiche d'assertion
(section~\ref{sec:c07-fiche}), les règles de qualité
(section~\ref{sec:c07-regles}), l'application au domaine~D
(section~\ref{sec:c07-lizbiz}) et les limites (section~\ref{sec:c07-limites}).

\separateur

% ============================================================================
\section{Pourquoi Décomposer}
\label{sec:c07-pourquoi}
% ============================================================================

\subsection{La borne de la mémoire de travail}

Un auditeur ne peut pas examiner « le contrôle d'accès » d'un bloc. La
psychologie cognitive établit depuis les travaux de George Miller que la
mémoire de travail ne retient qu'un petit nombre d'éléments à la fois,
de l'ordre de sept, plus ou moins deux. L'auditologie en fait l'un de ses
axiomes (axiome~2, \emph{cognition limitée}) et en déduit une
proposition~: tout instrument qui couvre un espace d'audit non trivial doit
être organisé en une décomposition hiérarchique, dont chaque niveau respecte
cette borne~\cite{minka2026auditologie}.

\begin{boxdef}
\termecle{Espace d'audit}\index{espace d'audit}~: ensemble des objets sur
lesquels porte une mission (processus, contrôles, données, traces), tel que
le délimite la lettre de mission.

\medskip
\termecle{Décomposition}\index{décomposition}~: découpage hiérarchique de
l'espace d'audit en niveaux emboîtés, chaque élément d'un niveau se divisant
en éléments du niveau inférieur.
\end{boxdef}

\subsection{Ce que la décomposition apporte}

\begin{enumerate}
    \item \textbf{La couverture}~: on sait ce qui a été audité, et ce qui ne
          l'a pas été.
    \item \textbf{La traçabilité}~: chaque conclusion remonte jusqu'à une
          trace dans le système, ce qui la rend opposable (axiome~6).
    \item \textbf{La localisation de l'incertitude}~: les réponses \NEK{} se
          situent à un endroit précis de l'arbre, et non dans « le domaine »
          (section~\ref{sec:c07-lizbiz}).
    \item \textbf{La répartition du travail}~: chaque sous-domaine peut être
          confié à un auditeur.
\end{enumerate}

% ============================================================================
\section{Le Gabarit Canonique}
\label{sec:c07-gabarit}
% ============================================================================

\subsection{Huit niveaux}

Le programme de l'auditologie propose un \termecle{gabarit
canonique}\index{gabarit canonique} de décomposition à huit
niveaux~\cite{minka2026auditologie}. Sa justification formelle (borne
cognitive, exhaustivité, non-redondance) fait l'objet d'un article en cours
de formalisation.

\noindent\begin{minipage}{\linewidth}
\begin{center}
\begin{tikzpicture}[
    niv/.style = {
        draw=CouleurBordDef, fill=CouleurDef, rounded corners=3pt,
        minimum width=3.1cm, minimum height=0.62cm, align=center,
        font=\footnotesize\bfseries\color{CouleurPrimaire}
    },
    note/.style = {font=\scriptsize\itshape\color{CouleurSecondaire}, anchor=west, text width=6.4cm},
    fleche/.style = {-{Stealth[length=4pt]}, thick, color=CouleurSecondaire}
]
\foreach \t/\q [count=\i] in {
    1. Domaine/Que soumet-on à l'audit~?,
    2. Sous-domaine/Quels grands ensembles le composent~?,
    3. Processus/Quelles suites d'activités~?,
    4. Activité/Quelles actions élémentaires~?,
    5. Contrôle/Quel dispositif maîtrise l'activité~?,
    6. Test/Comment l'auditeur le vérifie-t-il~?,
    7. Preuve/Quel élément le test produit-il~?,
    8. Trace/Quel enregistrement du système l'atteste~?} {
    \node[niv] (n\i) at (0, -0.85*\i) {\t};
    \node[note] at (1.8, -0.85*\i) {\q};
}
\foreach \i [evaluate=\i as \j using int(\i+1)] in {1,...,7}
    \draw[fleche] (n\i) -- (n\j);
\draw[decorate, decoration={brace, amplitude=5pt, mirror}, thick, color=CouleurAccent]
    (-1.75,-0.55) -- (-1.75,-3.65) node[midway, left=6pt, font=\scriptsize, align=right,
    color=CouleurAccent] {ce qui\\ est audité};
\draw[decorate, decoration={brace, amplitude=5pt, mirror}, thick, color=CouleurAccent]
    (-1.75,-3.95) -- (-1.75,-7.05) node[midway, left=6pt, font=\scriptsize, align=right,
    color=CouleurAccent] {comment on\\ le vérifie};
\end{tikzpicture}
\end{center}
\vspace{-0.4em}
\begin{center}{\small\itshape Figure~7.1. Le gabarit canonique de
décomposition (d'après~\cite{minka2026auditologie})}\end{center}
\end{minipage}

\medskip

Les quatre premiers niveaux décrivent \emph{ce qui est audité}~: ils
découpent le système. Les quatre derniers décrivent \emph{comment on le
vérifie}~: du dispositif de maîtrise jusqu'à l'enregistrement qui en prouve
le fonctionnement. Le contrôle fait la charnière~: c'est l'objet sur lequel
l'auditeur se prononce.

\subsection{Les huit niveaux en détail}

\begin{tableaucours}{L{2.3cm}L{5.0cm}L{4.9cm}}{%
    \tableauHeader{Niveau} &
    \tableauHeader{Définition} &
    \tableauHeader{Exemple (domaine D de LiZbiZ)}
}
Domaine & Champ d'audit défini par la lettre de mission & Contrôle d'accès
à l'application intégrée \\
Sous-domaine & Grand ensemble cohérent du domaine & Cycle de vie des
comptes \\
Processus & Suite d'activités qui produit un résultat & Départ d'un
agent \\
Activité & Action élémentaire du processus & Retrait des accès \\
Contrôle & Dispositif qui maîtrise l'activité & Désactivation du compte
sous 48~heures \\
Test & Procédure par laquelle l'auditeur vérifie le contrôle & Rapprocher
les départs et les comptes actifs \\
Preuve & Élément produit et conservé par le test & Extraction des comptes
actifs, avec son empreinte \\
Trace & Enregistrement du système qui atteste le fonctionnement &
Entrée du journal de désactivation \\
\end{tableaucours}
\vspace{-0.8em}
\begin{center}{\small\itshape Tableau~7.1. Les niveaux du gabarit
canonique}\end{center}

\begin{boxalert}
\textbf{Violation typique.} Placer une politique (« une politique d'accès
approuvée existe ») au niveau du test, et un journal technique au niveau du
sous-domaine, dans la même liste.

\smallskip
\textbf{Règle~:} chaque élément se place à un seul niveau, celui qui
correspond à sa nature. Un mélange de niveaux rend impossible le contrôle de
la couverture.
\end{boxalert}

% ============================================================================
\section{La Fiche d'Assertion}
\label{sec:c07-fiche}
% ============================================================================

\subsection{Cinq attributs par contrôle}

Le gabarit dit \emph{où} se trouve un contrôle. La \termecle{fiche
d'assertion}\index{fiche d'assertion} dit \emph{ce que l'auditeur affirme}
à son sujet, et ce qui arrive si l'affirmation est fausse.

\begin{tableaucours}{L{2.4cm}L{5.0cm}L{4.8cm}}{%
    \tableauHeader{Attribut} &
    \tableauHeader{Rôle} &
    \tableauHeader{Forme}
}
Objectif & Ce que le contrôle doit garantir & À l'infinitif~: « S'assurer
que\ldots » \\
Critère & La condition vérifiable qui définit la conformité & Au présent~:
une condition unique, observable \\
Question & La question posée pour coter le contrôle & Interrogative~;
réponse \OUI{}, \NON{} ou \NEK{} \\
Risque & L'état de non-conformité si la réponse est \NON{} & L'antithèse
du critère \\
Conséquence & L'effet sur l'organisation si le risque se réalise &
Financier, opérationnel, réglementaire \\
\end{tableaucours}
\vspace{-0.8em}
\begin{center}{\small\itshape Tableau~7.2. Les attributs de la fiche
d'assertion}\end{center}

La fiche est complétée par les niveaux inférieurs du gabarit~: le test,
la preuve attendue et la trace. Le prologue en a donné un premier exemple
(contrôle D.2.1)~; en voici un second.

\begin{boxobjet}{D.4.2 · Revue des comptes à privilèges}
\textbf{Position~:} domaine D $\rightarrow$ sous-domaine D.4 (comptes à
privilèges) $\rightarrow$ processus «~gestion des comptes d'administration~»
$\rightarrow$ activité «~revue des comptes~» $\rightarrow$ contrôle D.4.2.

\textbf{Objectif~:} s'assurer que chaque compte d'administration de
l'application correspond à un besoin en cours.

\textbf{Critère~1~:} la liste des comptes d'administration est revue chaque
trimestre par le responsable de la sécurité.

\textbf{Question~1~:} la revue des comptes d'administration a-t-elle été
réalisée pour chaque trimestre de la période~?

\textbf{Critère~2~:} le compte rendu de chaque revue est signé par le
responsable de la sécurité.

\textbf{Question~2~:} le compte rendu de chaque revue trimestrielle est-il
signé par le responsable de la sécurité~?

\textbf{Test~:} demander les comptes rendus des deux trimestres de la
période~; comparer leur liste à l'extraction des comptes d'administration.

\textbf{Preuve et trace~:} comptes rendus signés~; extraction des comptes
d'administration, avec son empreinte~; dates de dernière connexion de ces
comptes dans le journal.

\textbf{Risque~:} un compte d'administration inutile reste actif et sert à
modifier des données ou des droits sans contrôle.

\textbf{Conséquence~:} fraude ou altération non détectée, perte de fiabilité
des états de stock, manquement à l'obligation de sécuriser le système
(loi n°~2010/012, art.~26).
\end{boxobjet}

\subsection{De la fiche à la cotation}

La question de la fiche est celle que l'auditeur cote \OUI{}, \NON{} ou
\NEK{} (Chapitre~\ref{chap:preuve}). Le risque est \emph{activé} lorsque la
réponse est \NON{}~; il reste \emph{indéterminé} lorsque la réponse est
\NEK{}, et c'est précisément ce que mesure le taux d'incertitude~$U$.

Un contrôle peut porter plusieurs critères, comme D.4.2~: une revue
réalisée mais non signée ne satisfait pas le contrôle, et l'auditeur doit
pouvoir dire lequel des deux fait défaut. Chaque critère reçoit sa propre
question et sa propre cote. La cote du contrôle s'en déduit~: un seul
\NON{} suffit pour \NON{}~; à défaut, un seul \NEK{} donne \NEK{}~; le
contrôle n'est coté \OUI{} que si tous ses critères le sont.

% ============================================================================
\section{Les Règles de Qualité}
\label{sec:c07-regles}
% ============================================================================

\subsection{L'atomicité}

\begin{boxinfo}
\textbf{Règle d'atomicité.} L'objectif, le critère et la question ne
contiennent aucun connecteur logique «~et~» ni «~ou~». Toute assertion qui
en contient est scindée en autant d'assertions qu'il y a de conditions
indépendantes.
\end{boxinfo}

L'atomicité garantit qu'une réponse \NON{} désigne un seul défaut. Une
question double, «~les comptes des agents partis sont-ils désactivés
\emph{et} leurs données archivées~?~», n'admet pas de réponse honnête
lorsque les comptes sont désactivés mais les données non archivées~: le
répondant choisit au hasard, ou répond \NEK{} à tort.

\begin{boxalert}
\textbf{Violation typique.} « Les mots de passe sont-ils robustes et
changés régulièrement~? »

\smallskip
\textbf{Règle~:} scinder~: « La politique impose-t-elle au moins douze
caractères~? » et « Les mots de passe des comptes d'administration ont-ils
été changés au cours des 90 derniers jours~? ». Chaque question a son
critère, sa preuve et sa trace.
\end{boxalert}

\subsection{Exhaustivité, non-redondance, borne cognitive}

Une décomposition est de qualité lorsqu'elle respecte trois conditions,
que le programme de recherche formalise (article en préparation).

\begin{tableaucours}{L{3.2cm}L{5.0cm}L{4.0cm}}{%
    \tableauHeader{Condition} &
    \tableauHeader{Sens} &
    \tableauHeader{Défaut si elle manque}
}
Exhaustivité & Les éléments d'un niveau couvrent entièrement leur parent &
Une partie du domaine n'est pas auditée \\
Non-redondance & Deux éléments d'un même niveau ne se recouvrent pas~;
chaque élément dérive d'un seul parent & Double comptage, conclusions
contradictoires \\
Borne cognitive & Aucun élément n'a plus de sept enfants environ & Grille
illisible, éléments oubliés \\
\end{tableaucours}
\vspace{-0.8em}
\begin{center}{\small\itshape Tableau~7.3. Les trois conditions d'une bonne
décomposition}\end{center}

\begin{boxprocedure}
\textbf{La décomposition dans le classeur O4.} Pour que la décomposition
puisse être revue par le chef de mission, le classeur indique, pour chaque
élément~:
\begin{itemize}
    \item son code (par exemple D.4.2), son niveau et son parent~;
    \item pour chaque contrôle, sa fiche d'assertion complète~;
    \item pour chaque test, la preuve attendue et la référence du papier de
          travail où elle sera classée~;
    \item le référentiel d'où provient le critère (procédure interne,
          COBIT, ISO/IEC~27002, loi)~;
    \item la date de la revue de la décomposition et le nom du réviseur.
\end{itemize}
\end{boxprocedure}

\subsection{Relier la décomposition aux référentiels}

Les référentiels de contrôle fournissent des critères~: la
mesure~5.18 de l'ISO/IEC~27002:2022 sur les droits d'accès, les pratiques de
l'objectif DSS05 de COBIT, les principes du COSO
(Chapitre~\ref{chap:referentiels}). Ils ne fournissent pas la
décomposition de \emph{votre} mission~: c'est l'auditeur qui la construit, à
partir de la lettre de mission et de la prise de connaissance, puis va
chercher dans les référentiels les critères de chaque contrôle.

% ============================================================================
\section{Application : le Domaine D de LiZbiZ}
\label{sec:c07-lizbiz}
% ============================================================================

\subsection{Cinq sous-domaines}

\begin{boxscenario}{Opération LIZBIZ. Décomposer le domaine D}
AUDIX CONSEIL décompose le domaine~D (contrôle d'accès à l'application
intégrée) en cinq sous-domaines, eux-mêmes découpés en cinq éléments au
plus, soit $5 \times 5$ éléments.

\medskip
\begin{tableaucours}{cL{3.6cm}L{6.8cm}}{%
    \tableauHeader{Code} &
    \tableauHeader{Sous-domaine} &
    \tableauHeader{Processus}
}
D.1 & Politique et organisation & Politique d'accès~; rôles~; profils
types~; validation des demandes~; communication \\
D.2 & Cycle de vie des comptes & Arrivée~; changement de fonction~; départ~;
comptes génériques~; revue périodique \\
D.3 & Authentification & Mots de passe~; double authentification~;
verrouillage~; sessions~; comptes par défaut \\
D.4 & Comptes à privilèges & Attribution~; revue~; usage tracé~; comptes de
service~; secours \\
D.5 & Journalisation et revue & Journaux d'accès~; protection des
journaux~; revue des anomalies~; conservation~; alertes \\
\end{tableaucours}

\medskip
\textbf{Vérification.} Chaque sous-domaine compte cinq processus, sous la
borne de sept. Les cinq sous-domaines couvrent l'ensemble du contrôle
d'accès décrit dans la lettre de mission. Le compte générique de
l'agence de Garoua (Chapitre~\ref{chap:controle-interne}) relève de D.2 et
non de D.4~: il n'a pas de privilèges d'administration.
\end{boxscenario}

\subsection{Localiser l'incertitude}

\begin{boxlecture}{Domaine D décomposé}
Le prologue a lu le domaine~D d'un bloc~: 20 réponses, $\mu = 0{,}65$,
$U = 0{,}25$, $H \approx 1{,}24$~bit. La décomposition permet de répartir ces
réponses~: une question par sous-domaine, posée aux quatre mêmes
interlocuteurs.

\medskip
\begin{tableaucours}{cL{3.6cm}cccc}{%
    \tableauHeader{Code} &
    \tableauHeader{Sous-domaine} &
    \tableauHeader{\OUI{}/\NON{}/\NEK{}} &
    \tableauHeader{$\mu$} & \tableauHeader{$U$} & \tableauHeader{$H$}
}
D.1 & Politique et organisation & 4 / 0 / 0 & 1,00 & 0,00 & 0,00 \\
D.2 & Cycle de vie des comptes & 1 / 1 / 2 & 0,25 & 0,50 & 1,50 \\
D.3 & Authentification & 4 / 0 / 0 & 1,00 & 0,00 & 0,00 \\
D.4 & Comptes à privilèges & 2 / 0 / 2 & 0,50 & 0,50 & 1,00 \\
D.5 & Journalisation et revue & 2 / 1 / 1 & 0,50 & 0,25 & 1,50 \\
\end{tableaucours}

\medskip
\textbf{Constatation.} Toute l'incertitude du domaine se concentre dans
trois sous-domaines~: D.2, D.4 et D.5. Les sous-domaines D.1 et D.3 sont
connus et conformes.

\medskip
\textbf{Analyse} (confiance moyenne). Lu d'un bloc, le domaine paraissait
uniformément opaque. Décomposé, il montre que la politique et
l'authentification sont maîtrisées, et que l'ignorance porte sur la vie des
comptes, les privilèges et les journaux~: c'est là que les tests doivent
porter. D.1 et D.3, avec une entropie nulle et $U = 0$, sont des
certitudes~; à l'inverse, une entropie nulle avec $U = 1$ serait une
ignorance unanime (Chapitre~\ref{chap:conclure}).

\medskip
\hfill\lecture{0{,}65}{0{,}25}{1{,}24}
\end{boxlecture}

\begin{boxmetier}{Auditeur SI}
\textbf{Construire la décomposition d'un domaine}
\begin{enumerate}
    \item Partir du domaine défini par la lettre de mission, sans
          l'élargir.
    \item Découper en cinq à sept sous-domaines, à partir de la prise de
          connaissance et des référentiels.
    \item Descendre jusqu'aux contrôles, puis rédiger leurs fiches
          d'assertion.
    \item Pour chaque contrôle, définir le test, la preuve attendue et la
          trace à rechercher.
    \item Faire revoir la décomposition par le chef de mission avant le
          terrain, en vérifiant les trois conditions.
\end{enumerate}
\end{boxmetier}

% ============================================================================
\section{Les Limites du Gabarit}
\label{sec:c07-limites}
% ============================================================================

Le gabarit n'est pas une loi universelle. Il suppose trois conditions sur le
domaine audité, dont l'échec appelle une adaptation, et non l'abandon du
gabarit.

\begin{tableaucours}{L{3.6cm}L{4.4cm}L{4.2cm}}{%
    \tableauHeader{Condition} &
    \tableauHeader{Situation où elle manque} &
    \tableauHeader{Adaptation}
}
Richesse suffisante & Un audit de conformité qui vérifie seulement
l'existence de documents & Gabarit tronqué (domaine, sous-domaine,
contrôle, test) \\
Cohérence sémantique & Un domaine nouveau, sans regroupements stables
(audit d'un système d'IA) & Décomposition provisoire, qui converge vers le
gabarit \\
Hétérogénéité bornée & Une mission qui mêle sécurité, finance et risques
opérationnels & Plusieurs arbres, un par domaine \\
\end{tableaucours}
\vspace{-0.8em}
\begin{center}{\small\itshape Tableau~7.4. Conditions de validité et
adaptations du gabarit}\end{center}

\begin{boiteRemarque}{Et la première version de la méthode~?}
Une première version de la méthode du cours décomposait l'assertion
d'audit elle-même en huit niveaux~: domaine, point de contrôle, objectif,
critère, élément requis, question, risque, conséquence. Ces notions sont
conservées, mais à leur place~: les deux premières correspondent aux
niveaux supérieurs du gabarit, l'élément requis devient la preuve, et les
cinq autres forment la fiche d'assertion de chaque contrôle. Le gabarit
découpe le \emph{système}~; la fiche porte l'\emph{affirmation}.
\end{boiteRemarque}

\separateur

\begin{boxresume}
\textbf{Ce qu'il faut retenir de ce chapitre~:}
\begin{itemize}
    \item \textbf{On décompose parce que la mémoire de travail est
          limitée}~: environ sept éléments par niveau.
    \item \textbf{Le gabarit canonique} a huit niveaux~: domaine,
          sous-domaine, processus, activité (ce qui est audité), contrôle,
          test, preuve, trace (comment on le vérifie).
    \item \textbf{La fiche d'assertion} donne à chaque contrôle un objectif,
          un critère, une question, un risque et une conséquence.
    \item \textbf{L'atomicité} interdit «~et~» et «~ou~» dans l'objectif, le
          critère et la question.
    \item \textbf{Une bonne décomposition} est exhaustive, non redondante,
          et respecte la borne cognitive.
    \item \textbf{La décomposition localise l'incertitude}~: elle montre où
          porte l'ignorance, et donc où tester.
\end{itemize}
\end{boxresume}

\bigskip

\begin{boxexo}[title={\ding{45}\quad Exercice~7.1. À quel niveau~? \niveaubase}]
Placez chaque élément dans le niveau du gabarit qui lui correspond.
\begin{enumerate}[(a)]
    \item La revue trimestrielle des droits par le responsable de la
          sécurité.
    \item Une ligne du journal indiquant la désactivation du compte
          « afouda » le 3~mai.
    \item La gestion des sauvegardes.
    \item L'extraction des comptes actifs au 30~avril, avec son empreinte.
    \item Rapprocher la liste des départs et celle des comptes actifs.
    \item Le départ d'un agent.
\end{enumerate}
\end{boxexo}

\begin{boxexo}[title={\ding{45}\quad Exercice~7.2. Rendre atomique \niveaumoyen}]
Réécrivez les assertions suivantes pour qu'elles respectent la règle
d'atomicité, et donnez pour chacune le critère et la preuve attendue.
\begin{enumerate}[(a)]
    \item « Les sauvegardes sont réalisées chaque nuit et testées chaque
          mois. »
    \item « Les comptes génériques sont interdits ou placés sous la
          responsabilité d'un titulaire nommé. »
    \item « S'assurer que les journaux sont conservés et protégés contre la
          modification. »
\end{enumerate}
\end{boxexo}

\begin{boxexo}[title={\ding{45}\quad Exercice~7.3. Décomposer le domaine des stocks \niveauexpert}]
La seconde zone d'audit de LiZbiZ est la gestion des stocks dans
l'application intégrée.
\begin{enumerate}[(a)]
    \item Proposez cinq sous-domaines et, pour l'un d'eux, ses processus et
          ses activités.
    \item Rédigez la fiche d'assertion complète de deux contrôles, dont celui
          qui porte sur la validation des ajustements de stock.
    \item Vérifiez les trois conditions de qualité de votre décomposition.
    \item Expliquez où vous placeriez le fichier tableur de suivi des stocks
          tenu à Garoua.
\end{enumerate}
\end{boxexo}

\bigskip

\begin{boxinfo}
\textbf{Transition.} La décomposition dit quoi vérifier et comment. Il
reste à aller chercher les réponses et les preuves sur le terrain. Le
Chapitre~\ref{chap:preuve} présente l'instrument \OUI{}~/ \NON{}~/ \NEK{},
les entretiens, les tests, l'échantillonnage et les papiers de travail.
\end{boxinfo}
